\begin{helper}[The abstract comparison scheme: Morrey, Campanato and H\"older]
Let $n\ge2$, $p>1$, $0<\tau<n$, $c_0>0$, $W\ge0$, $M\ge0$, $d>0$ and $R_1>0$,
let $\sigma_2\colon\R\to\R$ be nonnegative, continuous, increasing with
$\sigma_2(0)=0$, and suppose the Sobolev--Poincar\'e inequality on balls holds in
dimension $n$ (\noderef{CitedSobolevPoincareBall}).  Write
$q=\frac n{n-1}$ and $\theta=n+(1-\frac\tau p)\frac n{n-1}$.  Then there are
$N_1,N_2\ge0$, depending only on these quantities, with the following property.

Let $0\le\varepsilon\le1$, let $\Omega,\Omega'',\Omega'\subseteq\R^n$ be open
with $\Omega'$ nonempty and a bounded Lipschitz domain
(\noderef{IsBoundedLipschitzDomain}) --- which is what the cited Campanato
embedding (C5) assumed below now requires of it --- $\Omega''$ bounded,
$\overline{\Omega'}\subseteq\Omega''$ and $\overline{\Omega''}\subseteq\Omega$,
and with the two separations
\[
d\le|x-y|\quad\text{for }x\in\Omega',\ y\notin\Omega'' ,
\qquad
d\le|x-y|\quad\text{for }x\in\Omega'',\ y\notin\Omega .
\]
Assume the cited Campanato embedding (C5) on $\Omega'$ at the exponent $q$
(\noderef{CitedCampanatoEmbedding}).  Let $u\in W^{1,p}(\Omega)$ with weak
gradient $g=\nabla u$ (\noderef{MemW1pOn}) satisfy $|u|\le M$ almost everywhere
on $\Omega$ and
\[
\int_\Omega|u|^p\,dx+\int_\Omega|\nabla u|^p\,dx\le W^p .
\]
Assume finally that for every $x_0\in\Omega''$ and every $0<R\le R_1$ with
$\overline B_R(x_0)\subseteq\Omega$ there is a comparison map
$v\in W^{1,p}\bigl(B_R(x_0)\bigr)$ with weak gradient $g_v$
(\noderef{MemW1pOn}) such that
\begin{align}
\int_{B_\rho(x_0)}(\varepsilon+|g_v|^2)^{\frac p2}dx
&\le c_0\Bigl(\frac\rho R\Bigr)^{n}\int_{B_R(x_0)}(\varepsilon+|g_v|^2)^{\frac p2}dx
\qquad(0<\rho\le R),\tag{4.11}\\
\int_{B_R(x_0)}|\nabla u-g_v|^p\,dx
&\le\sigma_2(R)\int_{B_R(x_0)}(1+|\nabla u|^2)^{\frac p2}dx .\tag{$\star$}
\end{align}
Then
\begin{itemize}
\item[(i)] $\|\,|\nabla u|\,\|_{L^{p,n-\tau}(\Omega')}\le N_1$
  (\noderef{MorreyBoundOn});
\item[(ii)] $\|u\|_{\mathcal L^{q,\theta}(\Omega')}\le N_2$
  (\noderef{CampanatoBoundOn});
\item[(iii)] for each $0\le\eta<1$ there is a function $w$ with $u=w$ almost
  everywhere on $\Omega'$ and $w\in C^\eta(\overline{\Omega'})$
  (\noderef{MemHolderOn}).
\end{itemize}

\noindent\textbf{What this node is, and why it carries all three conclusions.}
The thesis proves Theorems 3 and 4 and then disposes of Theorems 5 and 6 in one
sentence: ``In proving Theorems 3 and 4 the only specific properties of $u$, $v$
and $w$ we needed were contained in Lemma 4.1 and \S3.  Therefore, since we have
Lemma 5.1, the same proofs of Theorems 3 and 4 give Theorems 5 and 6 below''
(p.~50).  This node is that sentence made formal: the comparison maps enter only
through the two displayed estimates (4.11) of the thesis and $(\star)$, and the
function $u$ only through its Sobolev membership, its $L^\infty$ bound and the
$W^{1,p}$ datum, so everything the two settings share --- the Morrey decay
(4.14) of the thesis obtained by iteration, its passage to the Morrey bound (i), the
passage from (i) to the Campanato bound (ii), and the passage from (ii) to the
H\"older representative (iii) --- is proved once, here.  What the two theorems
do \emph{not} share, and therefore keep, is exactly what the thesis' sentence
names: the construction of the comparison map $v$ from their own structure
conditions, (F1) of the thesis--(F3) of the thesis in one case and
(Q1) of the thesis--(Q6) of the thesis in the other, together with the frozen operator,
the continuous representative and the two constants $c_0$ and $\sigma_2$ that
Theorem~1 and Lemma~4.1 (respectively Lemma~5.1) supply.

The comparison bundle is \emph{$\tau$-free}: neither (4.11) of the thesis nor
$(\star)$ mentions $\tau$.  That is what allows (iii) to be proved here, since
its proof re-runs the argument of (i) and (ii) at a second exponent
$t\in(0,n)$ depending on $\eta$, with the same bundle.

\noindent\textbf{On the $W$ datum.}  The hypothesis bounds
$\int_\Omega|u|^p+\int_\Omega|\nabla u|^p$, i.e.\ the genuine Sobolev norm, and
not $\int_\Omega(1+|\nabla u|^2)^{p/2}dx$.  The latter integrand is pointwise at
least $1$, so a bound on it would force $W\ge|\Omega|$ and would be
unsatisfiable when $|\Omega|=\infty$; since the theorems of \S4 and \S5 that
build this bundle are stated for an arbitrary open $\Omega$, that would make
them vacuous on instances such as $\Omega=\R^n$, and it would also let the
constants degrade with $|\Omega|$.  The weighted integral is still available on
each ball, because $(1+t^2)^{p/2}\le2^{p/2}(1+t^p)$ for $t\ge0$ gives
\[
\int_{B_R}(1+|\nabla u|^2)^{\frac p2}dx
\;\le\;2^{\frac p2}\Bigl(|B_R|+\int_{B_R}|\nabla u|^p\,dx\Bigr)
\;\le\;2^{\frac p2}\bigl(\omega_nR^{\,n}+W^p\bigr),
\]
whose right-hand side involves only $n$, $p$, $R$ and $W$.
\end{helper}
\begin{proof}
Throughout, $B_t=B_t(x_0)$ and $\overline B_t=\overline B_t(x_0)$ for the ball
currently in play, $|E|$ is Lebesgue measure, $\omega_n=|B_1(0)|\in(0,\infty)$
(the Preamble item on the volume of a ball), and all non-integer powers are real
powers of nonnegative numbers.  We also write
\[
q=\frac n{n-1},\qquad
\theta=n+q\Bigl(1-\frac\tau p\Bigr)
=n+\Bigl(1-\frac\tau p\Bigr)\frac n{n-1},
\]
so that $q>1$ and $\theta$ is the Campanato exponent of conclusion (ii), and we
record once the only inequality about $\theta$ that is used:
\[
\theta>0 .
\tag{0}
\]
Indeed $n+q=\frac{n^2}{n-1}$ and $\frac{n+q}q=\frac nq+1=n$, so
$\theta=n+q-\frac{q\tau}p>0$ is equivalent to $\tau<p\,\frac{n+q}q=pn$, which
holds because $\tau<n$ and $p>1$.

The step numbers below are those of the thesis' proof.  Steps 2--5 there are
the passage (4.12)--(4.13)--(4.14) from the comparison bundle to the Morrey
bound for $\nabla u$; that passage has been extracted verbatim into
\noderef{MorreyBoundFromBundle} and is used here only through its conclusion,
so no step of that range recurs below.

\emph{Step 1: the constants, chosen before all the data.}  Apply
\noderef{MorreyBoundFromBundle} with the parameters $n,p,\tau,c_0,W,d,R_1$ and
the function $\sigma_2$.  Its hypotheses on those parameters are exactly the
present node's: $n\ge2$, $p>1$, $0<\tau<n$, $c_0>0$, $W\ge0$, $d>0$, $R_1>0$,
and $\sigma_2$ nonnegative, continuous, increasing with $\sigma_2(0)=0$.  Let
$N_1\ge0$ be the constant it produces; it depends only on
$n,p,\tau,c_0,W,d,R_1$ and $\sigma_2$, which is the quantifier order the
present statement prescribes for it.  Now put
$N_\flat=\max\bigl\{N_1,\ M\omega_n^{1/q}(d/4)^{n/q}\bigr\}$ and
  apply \noderef{MorreyGradientToCampanato} (Lemma~4.2) with the parameters
  $n,p,\tau$, with $N_\flat$ in the role of its $N$ and with $d/8$ in the role
  of its $d$: its hypotheses $2\le n$, $1<p$, $0<\tau\le n$, $0\le N_\flat$,
  $0<d/8$ hold, and its hypothesis \noderef{CitedSobolevPoincareBall} is the
  present node's.  Let $N'\ge0$ be the constant it produces; it depends only on
  $n,p,\tau,N_\flat$ and $d$.  Finally put
  \[
  N_2=2\max\Bigl\{N',\
  \bigl[M^{(q-p)^{+}}W^{\min\{p,q\}}(d/8)^{-\theta}\bigr]^{1/q},\
  \bigl[W^q\omega_n^{(1-q/p)^{+}}(d/8)^{\,n(1-q/p)^{+}-\theta}\bigr]^{1/q}\Bigr\}
  \ \ge0 ,
  \]
  where $t^{+}=\max\{t,0\}$.

Every constant named so far is built from $n,p,\tau,c_0,W,M,d,R_1$ and
$\sigma_2$ alone --- in particular not from
$\varepsilon,\Omega,\Omega'',\Omega',u,g$ --- which is the quantifier order the
statement prescribes.  Now fix those data satisfying all the hypotheses.

\emph{Step 6: part (i).}  Apply \noderef{MorreyBoundFromBundle} again, now to
the data just fixed, at the constant $N_1$ of Step 1.  Its remaining hypotheses
are the present node's: $0\le\varepsilon\le1$; $\Omega$ and $\Omega''$ open
with $\Omega''$ bounded and $\overline{\Omega''}\subseteq\Omega$; the
separation $d\le|x-y|$ for $x\in\Omega''$ and $y\notin\Omega$;
$u\in W^{1,p}(\Omega)$ with weak gradient $g$ (\noderef{MemW1pOn}); the
$W$ datum $\int_\Omega|u|^p+\int_\Omega|\nabla u|^p\le W^p$; and the comparison
bundle, which is the present node's final hypothesis verbatim --- the cited
node asks for exactly the pair (4.11) of the thesis, $(\star)$ at a comparison map
$v\in W^{1,p}(B_R(x_0))$ with weak gradient $g_v$, on the same range of
$(x_0,R)$.  Its conclusion is
\[
\|\,|g|\,\|_{L^{p,n-\tau}(\Omega'')}\le N_1
\qquad(\noderef{MorreyBoundOn}).
\tag{6}
\]

Part (i) follows by restriction: $|g|\in L^p(\Omega')$ because $\Omega'\subseteq\Omega''$
and $|g|\in L^p(\Omega'')$ by the $L^p$ clause of (6), and for
$x_1\in\Omega'\subseteq\Omega''$ and
$0<\rho<\operatorname{diam}\Omega'\le\operatorname{diam}\Omega''$
(the inequality because $\Omega'\subseteq\overline{\Omega'}\subseteq\Omega''$)
we get
$\int_{\Omega'\cap B_\rho(x_1)}|g|^p\le
\int_{\Omega''\cap B_\rho(x_1)}|g|^p\le N_1^p\rho^{\,n-\tau}$.
This is $\|\,|\nabla u|\,\|_{L^{p,n-\tau}(\Omega')}\le N_1$, i.e.\ (i).

\emph{Step 7: a mean-replacement inequality.}  Let $S\subseteq\R^n$ be
measurable with $0<|S|<\infty$, let $h\in L^q(S)$ be real-valued and let
$c\in\R$.  Writing $h_S=|S|^{-1}\int_Sh$ for the mean,
\[
\int_S|h-h_S|^q\,dx\ \le\ 2^q\int_S|h-c|^q\,dx .
\tag{M}
\]
This is \noderef{MeanReplacementLp} applied to the vector-valued function
$f=h\,e$, where $e\in\R^n$ is a fixed unit vector, at the exponent $q$ and at
the constant vector $c\,e$.  Its hypotheses hold: $q>1$ because $n\ge2$ gives
$q=\frac n{n-1}>1$; $S$ is measurable with $0<|S|<\infty$; and $f\in L^q(S)$
because $|f|=|h|$ pointwise and $h\in L^q(S)$.  Its conclusion transfers to
$(M)$ because the map $t\mapsto t\,e$ is linear and isometric, so
$\int_Sf=\bigl(\int_Sh\bigr)e$ and hence $f_S=h_S\,e$, while
$|f-f_S|=|h-h_S|$ and $|f-c\,e|=|h-c|$ pointwise.

\emph{Step 8: part (ii).}  We prove
$\|u\|_{\mathcal L^{q,\theta}(\Omega')}\le N_2$ (\noderef{CampanatoBoundOn}).
The $L^q$ clause holds: $|u|\le M$ almost everywhere on $\Omega$ and
$|\Omega'|<\infty$ --- the set $\Omega'$ is bounded, being contained in
$\overline{\Omega'}\subseteq\Omega''$ with $\Omega''$ bounded, and a bounded
measurable set has finite Lebesgue measure --- so
$\int_{\Omega'}|u|^q\le M^q|\Omega'|<\infty$ by monotonicity of the integral
(Preamble), and $u$ is almost everywhere strongly measurable on $\Omega'$
because it is so on $\Omega$ by \noderef{MemW1pOn}.

Fix $x_1\in\Omega'$ and $0<\rho<\operatorname{diam}\Omega'$, write
$S=\Omega'\cap B_\rho(x_1)$ --- a set of positive measure, since $\Omega'$ is
open and $x_1\in\Omega'$, and of finite measure --- and let $u_S$ be the mean of
$u$ over $S$.  There are two regimes.

\emph{Small radii: $\rho\le d/8$.}  Apply \noderef{MorreyGradientToCampanato}
(Lemma~4.2) with the constants of Step 1 and with the data
$\Omega$, $B_{d/4}(x_1)$ in the role of its $\Omega'$ and $B_{d/8}(x_1)$ in the
role of its $\Omega''$, the function $u$ and the field $g$.  Its hypotheses
hold:
\begin{itemize}
\item $\Omega$, $B_{d/4}(x_1)$ and $B_{d/8}(x_1)$ are open, the middle one
  bounded; $\overline B_{d/4}(x_1)\subseteq\Omega''$, because a point $z$ with
  $|z-x_1|\le\frac d4$ and $z\notin\Omega''$ would contradict the hypothesis
  $d\le|x-y|$ for $x\in\Omega'$, $y\notin\Omega''$ taken at the pair
  $(x_1,z)$, since $x_1\in\Omega'$ and $\frac d4<d$; then
  $\Omega''\subseteq\overline{\Omega''}\subseteq\Omega$; and
  $\overline B_{d/8}(x_1)\subseteq B_{d/4}(x_1)$ because $\frac d8<\frac d4$;
\item the separation: if $y\in B_{d/8}(x_1)$ and $z\notin B_{d/4}(x_1)$ then
  $|y-z|\ge|z-x_1|-|x_1-y|>\frac d4-\frac d8=\frac d8$, which is the required
  lower bound $d/8$;
\item $u\in W^{1,p}_{\mathrm{loc}}(\Omega)$ with weak gradient $g$
  (\noderef{MemW1pLocOn}), by clause (3) of \noderef{WeakGradientRestrict}
  applied on $\Omega$ to the hypothesis $u\in W^{1,p}(\Omega)$;
\item $\|\,|g|\,\|_{L^{p,n-\tau}(B_{d/4}(x_1))}\le N_\flat$: the $L^p$ clause is
  the restriction of $|g|\in L^p(\Omega)$, and for $y\in B_{d/4}(x_1)$ and
  $0<\rho'<\operatorname{diam}B_{d/4}(x_1)=d/2$ we have
  $B_{d/4}(x_1)\subseteq\Omega''$ and
  $\rho'<d/2\le\operatorname{diam}\Omega''$, the latter because $\Omega''$
  contains the ball $B_{d/4}(x_1)$ of diameter $d/2$, so $(6)$ of Step 6
  gives $\int_{B_{d/4}(x_1)\cap B_{\rho'}(y)}|g|^p
  \le\int_{\Omega''\cap B_{\rho'}(y)}|g|^p\le N_1^p\rho'^{\,n-\tau}
  \le N_\flat^p\rho'^{\,n-\tau}$;
\item $\|u\|_{L^q(B_{d/4}(x_1))}\le N_\flat$ in the sense of the case
  $\lambda=0$ of \noderef{MorreyBoundOn}: $u\in L^q$ of that ball as in the
  first paragraph of this step, and for every centre and every $\rho'>0$,
  \[
  \int_{B_{d/4}(x_1)\cap B_{\rho'}(y)}|u|^q\,dx
  \le M^q\bigl|B_{d/4}(x_1)\bigr|=M^q\omega_n(d/4)^n\le N_\flat^{\,q},
  \]
  by $|u|\le M$ almost everywhere, monotonicity of the integral and the
  definition of $N_\flat$ in Step 1.
\end{itemize}
Its conclusion is $\|u\|_{\mathcal L^{q,\theta}(B_{d/8}(x_1))}\le N'$, whose
exponent $n+\frac n{n-1}(1-\frac\tau p)$ is the $\theta$ of the present proof.
Taking in it the centre $x_1\in B_{d/8}(x_1)$ and the radius
$\rho\le d/8<\operatorname{diam}B_{d/8}(x_1)=d/4$, and noting
$B_{d/8}(x_1)\cap B_\rho(x_1)=B_\rho(x_1)$, we obtain
\[
\int_{B_\rho(x_1)}\bigl|u-u_{B_\rho(x_1)}\bigr|^q\,dx\le N'^{\,q}\rho^{\,\theta},
\]
where $u_{B_\rho(x_1)}$ is the mean of $u$ over $B_\rho(x_1)$.  Now $(M)$ with
$S=\Omega'\cap B_\rho(x_1)$ and $c=u_{B_\rho(x_1)}$, followed by
$S\subseteq B_\rho(x_1)$ and monotonicity of the integral, gives
\[
\int_S|u-u_S|^q\,dx\le2^q\int_S\bigl|u-u_{B_\rho(x_1)}\bigr|^q\,dx
\le2^q\int_{B_\rho(x_1)}\bigl|u-u_{B_\rho(x_1)}\bigr|^q\,dx
\le2^qN'^{\,q}\rho^{\,\theta}\le N_2^{\,q}\rho^{\,\theta},
\]
by the definition of $N_2$ in Step 1.

\emph{Large radii: $\rho>d/8$.}  Here we use $(M)$ with $c=0$:
$\int_S|u-u_S|^q\le2^q\int_S|u|^q$, and we bound $\int_S|u|^q$ in the two
possible regimes of the exponents.
\begin{itemize}
\item If $q\ge p$: $|u|^q=|u|^p|u|^{q-p}\le M^{q-p}|u|^p$ almost everywhere on
  $\Omega$, so $\int_S|u|^q\le M^{q-p}\int_\Omega|u|^p\le M^{q-p}W^p$ by
  monotonicity of the integral and the hypothesis
  $\int_\Omega|u|^p+\int_\Omega|g|^p\le W^p$.  Since $\theta>0$ by $(0)$ and
  $\rho>d/8$, we get $1\le(8\rho/d)^{\theta}$, hence
  \[
  \int_S|u-u_S|^q\le2^qM^{q-p}W^p\le2^qM^{q-p}W^p(d/8)^{-\theta}\rho^{\,\theta}
  \le N_2^{\,q}\rho^{\,\theta}.
  \]
\item If $q<p$: by H\"older's inequality on the set $S$ of finite measure, with
  the exponents $p/q>1$ and its conjugate (Preamble),
  \[
  \int_S|u|^q\,dx\le\Bigl(\int_S|u|^p\,dx\Bigr)^{q/p}|S|^{1-q/p}
  \le W^q\bigl(\omega_n\rho^n\bigr)^{1-q/p},
  \]
  using $|S|\le|B_\rho(x_1)|=\omega_n\rho^n$.  The exponent of $\rho$ on the
  right is $n(1-q/p)$, and
  \[
  n\Bigl(1-\frac qp\Bigr)-\theta
  =n-\frac{nq}p-n-q+\frac{q\tau}p
  =-q\Bigl(1+\frac{n}{p}-\frac\tau p\Bigr)
  =-\frac qp\bigl(p+n-\tau\bigr)<0 ,
  \]
  because $\tau<n<p+n$.  Hence $\rho^{n(1-q/p)-\theta}\le(d/8)^{n(1-q/p)-\theta}$
  for $\rho>d/8$, and
  \[
  \int_S|u-u_S|^q\le2^qW^q\omega_n^{1-q/p}(d/8)^{\,n(1-q/p)-\theta}\rho^{\,\theta}
  \le N_2^{\,q}\rho^{\,\theta}
  \]
  by the definition of $N_2$ in Step 1.
\end{itemize}
In both regimes and for all admissible $(x_1,\rho)$ the Campanato bound with
constant $N_2$ holds, which is (ii).

\emph{Step 9: part (iii).}  Let $0\le\eta<1$ and put
\[
t=\frac{\min\{n,p\}(1-\eta)}2 .
\]
Then $0<t<\min\{n,p\}$, so $t\in(0,n)$ --- indeed
$0<t\le\min\{n,p\}/2\le n/2<n$, with room to spare, the worst case being the
included endpoint $\eta=0$ --- and $t<p$.  Steps 1--8 used the parameter $\tau$
only through the hypothesis $0<\tau<n$: it enters
\noderef{MorreyBoundFromBundle}, whose $\tau$ is a parameter bound before its
constant $N_1$ and before all the data, the constants $N_\flat,N',N_2$ of
Step 1, and the exponents $n-\tau$ and $\theta$ of the conclusions, and
nowhere else.  In particular the comparison bundle, which is the one hypothesis
that is not a bound on $u$ or a property of the three sets, does not mention
$\tau$ at all, and the data
$\varepsilon,\Omega,\Omega'',\Omega',u,g$ and that bundle are now fixed.
Running Steps 1--8 again with $t$ in place of $\tau$, the same data and the same
bundle, therefore produces a constant $N_2(t)\ge0$ with
\[
\|u\|_{\mathcal L^{q,\theta(t)}(\Omega')}\le N_2(t),
\qquad \theta(t)=n+q\Bigl(1-\frac tp\Bigr).
\]
Because $0<t<p$ we have $0<1-\frac tp<1$ and hence
\[
n<\theta(t)<n+q .
\]
Put $N=\max\{N_2(t),N_3\}$, where $N_3\ge0$ is chosen so that
$\|u\|_{L^q(\Omega')}\le N_3$ in the sense of the case $\lambda=0$ of
\noderef{MorreyBoundOn}; such an $N_3$ exists, for instance
$N_3=M\bigl|\Omega'\bigr|^{1/q}$, since $|u|\le M$ almost everywhere and
$\int_{\Omega'\cap B_{\rho}(y)}|u|^q\le M^q|\Omega'|$ for every centre and
radius, and $u\in L^q(\Omega')$ by Step 8.  (Here $N_3$ is allowed to depend on
$\Omega'$: part (iii) asserts no bound on any norm, only the existence of a
H\"older representative, so no constant of the conclusion depends on it.)
Enlarging a Morrey or Campanato bound preserves it, so both
$\|u\|_{\mathcal L^{q,\theta(t)}(\Omega')}\le N$ and
$\|u\|_{L^{q}(\Omega')}\le N$ hold.

Apply the cited Campanato embedding (C5)
$\mathrm{CitedCampanatoEmbedding}(n,q,\Omega')$
(\noderef{CitedCampanatoEmbedding}), a hypothesis of the present node.  Its
first antecedent is that $\Omega'$ be a bounded Lipschitz domain
(\noderef{IsBoundedLipschitzDomain}), which is also a hypothesis of the present
node; discharge it with that, and read the predicate at the
exponent $\theta(t)$ --- admissible because $n<\theta(t)<n+q$ --- and at the
constant $N\ge0$.  It produces a constant $N_h\ge0$ and, applied to $u$ with the
two bounds just recorded, a function $w_0$ with
\[
u=w_0\ \text{a.e. on }\Omega',
\qquad
w_0\in C^{\eta'}\bigl(\overline{\Omega'}\bigr)\ (\noderef{MemHolderOn}),
\qquad
|w_0(x)-w_0(y)|\le N_h|x-y|^{\eta'}
\]
for $x,y\in\overline{\Omega'}$, where
$\eta'=\frac{\theta(t)-n}q=1-\frac tp$.  By the choice of $t$,
\[
\eta'=1-\frac tp\ge1-\frac{1-\eta}2=\frac{1+\eta}2\ge\eta ,
\]
the last step because $\eta\le1$.

It remains to lower the exponent from $\eta'$ to $\eta$, which is where the
boundedness of $\Omega'$ enters.  Put
$\Delta=\operatorname{diam}\overline{\Omega'}<\infty$ and
$C_3=N_h\max\{1,\Delta^{\eta'-\eta}\}$.  For $x,y\in\overline{\Omega'}$ with
$x\ne y$ we have $0<|x-y|\le\Delta$ and therefore
\[
|w_0(x)-w_0(y)|\le N_h|x-y|^{\eta'}
=N_h|x-y|^{\eta}|x-y|^{\eta'-\eta}\le C_3|x-y|^{\eta},
\]
using $\eta'-\eta\ge0$ and the monotonicity of $t\mapsto t^{\eta'-\eta}$ on
$[0,\infty)$; and for $x=y$ both sides vanish when $\eta>0$, while for $\eta=0$
the right-hand side is $C_3\ge0$ and the left-hand side is $0$.  Together with
the continuity of $w_0$ on $\overline{\Omega'}$, which is the first clause of
\noderef{MemHolderOn} at the exponent $\eta'$ and does not involve the exponent,
this gives $w_0\in C^{\eta}(\overline{\Omega'})$ (\noderef{MemHolderOn}).  With
$u=w_0$ almost everywhere on $\Omega'$, the function $w=w_0$ witnesses (iii) for
the given $\eta$.  As $\eta\in[0,1)$ was arbitrary, (iii) holds, and the proof
is complete.
\end{proof}
